\section{The \sys{} Framework}
\label{sec:framework}

\subsection{Identification rubric}
\sys{} decomposes a DID study into eleven auditable dimensions: parallel
trends, no anticipation, staggered-timing handling, SUTVA/spillovers,
control-group construction, specification, inference, sample-period selection,
concurrent policies, placebo/robustness, and data measurement. Each dimension is
expressed as an \emph{assumption} (what must hold for identification), a testable
\emph{implication} (what should therefore be observable), and the
\emph{evidence} a credible paper would report to support it. For a carbon-pricing
or emissions-trading evaluation, for instance, the parallel-trends dimension
expects an event-study or pre-trend test with null pre-period coefficients, while
no-anticipation turns on the gap between a scheme's \emph{announcement} and
\emph{implementation} dates. The rubric is declarative, stored in a configuration
file (\texttt{identification\_dimensions.yaml}), so the audited dimensions are
explicit and inspectable rather than buried in code. A companion flaw taxonomy
(\S\ref{sec:evaluation}) maps each known identification threat one-to-one onto a
dimension, which both documents the rubric's coverage and supplies the
perturbations used for evaluation.

\subsection{Bounded LLM pipeline}
The audit runs a fixed, deterministic sequence, \emph{decomposition
$\rightarrow$ extraction $\rightarrow$ assessment $\rightarrow$ localization
$\rightarrow$ report}, shown in the top band of Figure~\ref{fig:arch}. The stage
order is written in code and never chosen by a model. \emph{Decomposition}
instantiates the rubric for the paper at hand; \emph{extraction} gathers, per
dimension, the in-paper evidence relevant to that assumption; \emph{assessment}
evaluates each assumption--implication--evidence chain and emits a per-dimension
risk judgement with a rationale and cited evidence; \emph{localization}
aggregates these judgements into a risk map ranking where support is weakest; and
\emph{report} renders a transparent, evidence-grounded document for a human
reviewer.

Agency is confined to the two stages where open-ended judgement actually
helps: extraction and assessment. Each invokes a bounded ReAct-style loop with
three constraints: (i)~a \emph{small fixed tool set} passed in per call: %
\texttt{evidence\_search} over the paper's text, figures and tables,
\texttt{figure\_parse} for structured figure/table content, and
\texttt{policy\_lookup} for policy metadata such as announcement versus
implementation dates; (ii)~a \emph{hard step budget} (\texttt{max\_steps}) so a
run cannot wander indefinitely; and (iii)~a \emph{logged trace} of every step,
tool call, and cited piece of evidence. All other stages are plain orchestration
with no agent loop. In the configuration evaluated in \S\ref{sec:results} the
step budget is minimal (\texttt{max\_steps}{=}1), and the LLM assessment path
resolves to a fixed per-dimension sequence (deterministic retrieval, one
relevance-gate call, one adequacy call). The gate labels each retrieved passage
\emph{high}, \emph{partial}, or \emph{none}: a dimension with no relevant
passage is scored \texttt{unknown} without an adequacy call, while one with only
partially relevant passages is still judged (Appendix~\ref{app:retrieval}). We
therefore call the design \emph{bounded} rather than agentic: the evaluated
system is structured, dimension-scoped prompting over retrieved evidence with an
explicit abstention path, not an open-ended agent, and for an audit we treat
this near-determinism as a feature.

\paragraph{Why bound the agency.}
A fully autonomous agent would make per-run behaviour path-dependent and
high-variance, undermining the central evaluation claim of
\S\ref{sec:evaluation}: that flaw injection provides \emph{reproducible} local
ground truth. Restricting agency to two clearly scoped sub-tasks preserves
reproducibility (stable, comparable metrics) and evidence-traceability (every
judgement anchored to logged evidence a human can inspect). The expert sits at
the end of the pipeline: \sys{} surfaces and localizes identification risk; the
judgement itself stays with the human.

\subsection{Output}
\sys{} emits study--dimension--risk triplets, each with an
evidence-grounded rationale, alongside the human-readable report. It never
produces a verdict on the true causal effect. This output shape keeps the
system both auditable by a human and measurable by the evaluation that follows.

\begin{figure*}[t]
  \centering
  \includegraphics[width=0.88\textwidth]{figures/figure2}
  \caption{\sys{} architecture. The top band is the audit pipeline, a fixed
  deterministic sequence; model calls are confined to two stages
  (extraction and assessment), each a bounded ReAct loop with a fixed tool set
  (evidence search, figure parse, policy lookup), a hard step budget, and a
  logged trace. The bottom band shows the flaw-injection evaluation loop, which
  injects one known flaw, audits the clean and injected versions, and scores
  detection, false alarm, and localization against the injected flaw as local
  ground truth, never against the unobserved true effect.}
  \label{fig:arch}
\end{figure*}
